\documentclass[12pt]{article}
\usepackage[letterpaper, margin=1in]{geometry}
\usepackage{setspace}
\doublespacing
\usepackage{amsmath,amssymb}
\usepackage{natbib}
\usepackage{hyperref}
\hypersetup{colorlinks=true,linkcolor=blue,citecolor=blue,urlcolor=blue}
\usepackage{microtype}

\setlength{\parindent}{0pt}
\setlength{\parskip}{6pt}

\begin{document}

% ============================================================
% TITLE PAGE
% ============================================================

\begin{center}
{\Large \textbf{Gravitational Decoherence on Timelike Worldlines:\\
The Missing Information Was Always Information}}

\vspace{12pt}

Heath W Mahaffey\\
Independent Researcher\\
2425 Entiat Way\\
Entiat, Washington, 98822
USA\\
\texttt{hmahaffeyges@gmail.com}

\vspace{12pt}

February 2026

\vspace{12pt}

\textit{Essay written for the Gravity Research Foundation\\
2026 Awards for Essays on Gravitation.}

\vspace{18pt}

\textbf{Abstract}
\end{center}

\begin{quote}
\noindent
From Aristotle to Galileo, and Newton to Einstein.  From uniform
acceleration to uniform gravitation.  Einstein revealed gravity as
geometry.  Jacobson revealed geometry as thermodynamics---the horizon
entropy functional determines the gravitational dynamics.  Each built upon
what came before.  Every application of Jacobson's result uses purely
geometric entropy.  Today we add another piece of the puzzle: the missing
term is classical information produced by gravitational decoherence of
matter on timelike worldlines.  Photons, propagating on null geodesics with
zero proper time, neither undergo nor produce decoherence.  Including
$S_{\mathrm{info}}$ in the entropy functional yields a dual-sector
gravitational coupling---$\mu(a) < 1$ for matter, $\Sigma(a) = 1$ for
photons---with zero free parameters.  The missing information was always
information.
\end{quote}

\newpage

% ============================================================
% ESSAY BODY
% ============================================================

\section*{I. The Missing Piece}

Jacobson showed us that gravity is thermodynamics---that the Einstein
field equations emerge as an equation of state from $\delta Q = T\,dS$
applied to causal horizons.  But he never identified the source of the
thermodynamics.  The entropy in his derivation was purely geometric:
\begin{equation}
S = \frac{A}{4\ell_P^2}.
\end{equation}
This applies to the horizon, but it doesn't explain the method.

Einstein faced a similar dilemma.  He added $\Lambda$ to his field
equations---a constant that balanced the geometry---but never identified
what it was.  It was a placeholder for something he could not yet name.

They were both missing the same piece of the puzzle.

$\Lambda$ is the vacuum baseline---the constant informational floor from
quantum fluctuations, always present, independent of structure.  The
missing piece is $S_{\mathrm{info}}$: the structural contribution, the
cumulative thermodynamic cost of gravitational decoherence of matter---irreversible
classical information encoded onto the cosmic horizon.
When matter collapses under gravity---into halos, galaxies, stars---it
forces quantum superpositions into definite classical states.  Each event
produces new classical information: a particle that was ``here or there''
becomes definitively ``here.''  That information, by Landauer's principle,
carries a thermodynamic cost.  By the holographic principle, it is encoded
on the horizon.

Einstein's cosmological constant and Jacobson's entropy functional were
both incomplete in the same way.  Neither included what gravity has been
producing all along.

This is not modified gravity.  The Einstein equations are untouched.  What
has changed is the entropy functional---the same place Jacobson showed the
field equations live.  Completing the entropy completes the gravitational
dynamics.


\section*{II. Duration Is the Key}

Why should this informational cost affect matter but not photons?  The
answer is proper time.

Massive particles propagate on timelike worldlines: $\Delta\tau > 0$.
Duration permits a temporal sequence of gravitational interactions.
Overdensities attract matter, matter collapses, quantum superpositions
resolve into classical outcomes.  Time is required for decoherence, and
only matter experiences time.

Photons propagate on null geodesics: $ds^2 = 0$, $\Delta\tau = 0$.  A
photon experiences no proper time.  It does not undergo gravitational
clustering.  It does not decohere.  It produces zero informational entropy.

This is not an assumption.  It is a consequence of the causal structure of
spacetime---the same causal structure that general relativity describes.
The metric signature distinguishes timelike from null.  This distinction
has thermodynamic consequences: matter pays the informational cost of
becoming classical, and photons do not.

When $S_{\mathrm{info}}$ is included in the entropy functional, its form
is derived, not chosen.  The coupling $\beta_m = \Omega_m/2$ follows from
the virial theorem---gravitationally bound systems in equilibrium partition
kinetic energy as half the potential energy.  The activation function
follows from horizon thermodynamics:
\begin{equation}
\mathcal{E}(a) = \exp\!\left(1 - \frac{1}{a}\right).
\end{equation}
At early times ($a \to 0$): $\mathcal{E} \to 0$.  No structure, no
decoherence, no cost.  At the present epoch ($a = 1$):
$\mathcal{E} = 1$.  Full activation.  In the far future
($a \to \infty$): $\mathcal{E} \to e$.  The feedback loop
self-regulates.  There will be no heat death, no empty dark and lonely
universe.

In the $\mu$--$\Sigma$ framework, the sector split maps precisely:
\begin{equation}
\mu(a) < 1, \qquad \Sigma(a) = 1.
\end{equation}
Matter feels a weakened effective gravitational coupling because part of the
energy budget is diverted to informational entropy production on the
horizon.  Photons, which produce no informational entropy, experience
standard gravity.  The coupling is derived from the virial theorem.  The
sector split is derived from the metric signature.  Zero free parameters.


\section*{III. The Downstream Effects}

If the missing piece fits, it should resolve existing puzzles without being
tuned to do so.  It does.

\textbf{The Hubble tension.}  The cosmic microwave background, observed
through the photon sector, yields
$H_0 = 67.4 \pm 0.5$~km\,s$^{-1}$\,Mpc$^{-1}$.  The local distance
ladder, built from matter-sector objects, yields
$H_0 = 73.0 \pm 1.0$~km\,s$^{-1}$\,Mpc$^{-1}$.  The $5\sigma$
discrepancy has resisted resolution for a decade.  With the informational
term included, the photon sector gives
$H_0^{\mathrm{photon}} = 67.16$~km\,s$^{-1}$\,Mpc$^{-1}$ and the matter
sector gives $H_0^{\mathrm{matter}} = 72.26$~km\,s$^{-1}$\,Mpc$^{-1}$---both
from a single MCMC posterior, both within $1\sigma$ of observations.  The
tension is not a discrepancy.  It is the empirical signature of the sector
split.

\textbf{The $S_8$ tension.}  The weakened gravitational coupling
($\mu < 1$) means matter perturbations grow less efficiently.  $\sigma_8$
drops from $0.809$ to $0.800$---in the direction that weak lensing
surveys (KiDS, DES, HSC) have been reporting for a decade.  The $S_8$
tension, like the Hubble tension, becomes a consequence rather than a
problem.

\textbf{Surgical precision.}  All standard cosmological parameters shift
by less than $0.1\sigma$.  The CMB power spectrum is virtually identical to
$\Lambda$CDM.  The modification targets only growth, because only growth
produces informational entropy.

\textbf{The arrow of time.}  Decoherence is irreversible---Landauer's
principle forbids costless erasure.  If gravity drives decoherence, and
decoherence produces the arrow of irreversibility, then the arrow of time
is gravitational in origin.

\textbf{The measurement problem.}  If gravitational decoherence on timelike
worldlines is the mechanism that produces classical reality, what does this
imply for quantum mechanics?  Einstein once remarked that God does not play
dice, and that spooky action at a distance did not feel right.  Was he
right?  Schr\"odinger's cat may never have been alive and dead---it was
always in a gravitationally decohering environment.  The superposition
would collapse not because a conscious observer opens the box, but because
gravity decoheres massive objects on timelike worldlines.  No conscious
observer required.  The double-slit experiment would tell the same story:
a photon passes through both slits because it propagates on a null
geodesic---$\Delta\tau = 0$, no decoherence, superposition preserved.
Place a massive detector at the slit and gravitational decoherence
collapses the which-path information.  The ``measurement'' was always
gravity.  If the same sector split that explains the Hubble
tension---$\mu < 1$ for matter, $\Sigma = 1$ for photons---also explains
why quantum mechanics works the way it does, the implications extend far
beyond cosmology.

\textbf{Self-regulation.}  Informational pressure drives expansion.
Expansion dilutes matter.  Diluted matter collapses less efficiently.
Reduced collapse produces less information.  $\mathcal{E}(a)$ asymptotes
rather than diverges.  The universe does not end in a Big Rip.  It
matures.

\textbf{The coincidence problem.}  Why does dark energy dominate at roughly
the same epoch as structure formation peaks?  Because the informational
term \textit{is driven by} structure formation.  The coincidence is not a
coincidence.  It is a consequence.

None of these results were tuned.  They emerge from a single addition to
the entropy functional with zero free parameters.


\section*{IV. Seventeen Chains}

These are not speculations.  The dual-sector mechanism has been implemented
in both MGCAMB and in direct Fortran modifications to CAMB~v1.5.8.
Seventeen MCMC chains have been completed against the full Planck~2018
likelihood.  All achieve Gelman--Rubin convergence $R - 1 < 0.01$.

The best-fit $\Delta\chi^2 = +0.54$ relative to $\Lambda$CDM---well within
the $95\%$ threshold of $3.84$.  The model is statistically
indistinguishable from $\Lambda$CDM in its fit to the CMB, while producing
the $\sigma_8$ suppression, the Hubble split, and the self-regulating
feedback as emergent consequences.

Zero free parameters beyond $\Lambda$CDM.  The coupling is derived.  The
activation function is derived.  The sector split is derived.  The model
either works with these values or it does not.  The chains say it works.

All chains, configuration files, modified source code, analysis scripts,
and the complete theoretical derivation are publicly archived at:\\
\url{https://github.com/hmahaffeyges/IAM-Validation}.


\section*{V. In the Path of Euclid}

Euclid is currently surveying 15,000~square degrees of sky to constrain the
$\mu$--$\Sigma$ parameter space.  Projected sensitivity:
$\sigma(\mu_0) \approx 0.04$.  The prediction: $\mu_0 = -0.135$.
Detection significance: $\sim 3.4\sigma$.

Most modified gravity theories predict $\mu \geq 1$.  The prediction
$\mu < 1$ with $\Sigma = 1$ occupies a region that is nearly empty
theoretically but fully accessible observationally.  Combined with
DESI Year~5, the projected significance reaches $4$--$5\sigma$.

The prediction is on the table awaiting the arrival of data.


\section*{VI. Filling the Missing Piece}

Jacobson's result was more powerful than anyone realized.  He showed that
the entropy functional determines the field equations---but was not able to
explain the source of the thermodynamics.  Einstein faced the same question
from the other direction: so he added $\Lambda$ to balance the geometry
while never being able to identify its source.

The informational term answers both.  $\Lambda$ is the vacuum
baseline---the constant informational floor.  $S_{\mathrm{info}}$ is the
structural contribution---the cost of gravity producing classical reality
from quantum raw material, one irreversible decoherence event at a time,
encoded on the cosmic horizon.

Gravity drives decoherence.  Decoherence produces information.  Information
modifies the horizon entropy.  The modified entropy changes the
gravitational dynamics.  The changed dynamics feed back into gravitational
collapse.  The loop closes.  Gravity is both the cause and the product.

Three equations.  Zero free parameters.  Seventeen converged chains.

The missing information was always information.


% ============================================================
% REFERENCES
% ============================================================

\newpage

\begin{thebibliography}{99}

\bibitem{Jacobson1995}
T.~Jacobson, ``Thermodynamics of Spacetime: The Einstein Equation of
State,'' \textit{Phys. Rev. Lett.} \textbf{75}, 1260 (1995).

\bibitem{CaiKim2005}
R.-G.~Cai and S.~P.~Kim, ``First law of thermodynamics and Friedmann
equations of Friedmann-Robertson-Walker universe,'' \textit{JHEP}
\textbf{02}, 050 (2005).

\bibitem{Bekenstein1973}
J.~D.~Bekenstein, ``Black holes and entropy,'' \textit{Phys. Rev. D}
\textbf{7}, 2333 (1973).

\bibitem{Zurek2003}
W.~H.~Zurek, ``Decoherence, einselection, and the quantum origins of the
classical,'' \textit{Rev. Mod. Phys.} \textbf{75}, 715 (2003).

\bibitem{Landauer1961}
R.~Landauer, ``Irreversibility and heat generation in the computing
process,'' \textit{IBM J. Res. Dev.} \textbf{5}, 183 (1961).

\bibitem{tHooft1993}
G.~'t~Hooft, ``Dimensional reduction in quantum gravity,'' in
\textit{Salamfestschrift} (World Scientific, 1993), gr-qc/9310026.

\bibitem{Susskind1995}
L.~Susskind, ``The world as a hologram,'' \textit{J. Math. Phys.}
\textbf{36}, 6377 (1995).

\bibitem{PogosianSilvestri2016}
L.~Pogosian and A.~Silvestri, ``What can cosmology tell us about gravity?
Constraining Horndeski gravity with $\Sigma$ and $\mu$,'' \textit{Phys.
Rev. D} \textbf{94}, 104014 (2016).

\bibitem{Planck2018}
Planck Collaboration, ``Planck 2018 results. VI. Cosmological parameters,''
\textit{Astron. Astrophys.} \textbf{641}, A6 (2020).

\bibitem{Riess2022}
A.~G.~Riess \textit{et al.}, ``A comprehensive measurement of the local
value of the Hubble constant,'' \textit{Astrophys. J. Lett.} \textbf{934},
L7 (2022).

\bibitem{Heymans2021}
C.~Heymans \textit{et al.} (KiDS), ``KiDS-1000 Cosmology,''
\textit{Astron. Astrophys.} \textbf{646}, A140 (2021).

\bibitem{DES2022}
DES Collaboration, ``Dark Energy Survey Year 3 results,'' \textit{Phys.
Rev. D} \textbf{105}, 023520 (2022).

\bibitem{Penrose1996}
R.~Penrose, ``On gravity's role in quantum state reduction,''
\textit{Gen. Relativ. Gravit.} \textbf{28}, 581 (1996).

\bibitem{Berut2012}
A.~B\'erut \textit{et al.}, ``Experimental verification of Landauer's
principle,'' \textit{Nature} \textbf{483}, 187 (2012).

\bibitem{Mahaffey2026obs}
H.~Mahaffey, ``Constraints on Late-Time $f\sigma_8$ Suppression from
$\mu < 1$, $\Sigma = 1$: Planck~2018 and Large-Scale Structure,''
submitted to \textit{Universe} (2026).

\end{thebibliography}

\end{document}
